\documentclass{article}
\usepackage{graphicx} % Required for inserting images
\usepackage[margin=1in]{geometry}
\usepackage{xcolor}
\usepackage{listings}

\definecolor{codebg}{rgb}{0.95,0.95,0.95}
\definecolor{coderule}{rgb}{0.8,0.8,0.8}

\lstdefinestyle{codebox}{
  backgroundcolor=\color{codebg},
  rulecolor=\color{coderule},
  frame=single,
  framerule=0.4pt,
  basicstyle=\ttfamily\small,
  breaklines=true,
  breakatwhitespace=false,
  columns=fullflexible,
  keepspaces=true,
  showstringspaces=false,
  xleftmargin=4pt,
  xrightmargin=4pt,
  framexleftmargin=4pt,
  framexrightmargin=4pt,
}
\lstset{style=codebox}

\title{WMUX - User Manual}
\author{Peter Kure}
\date{July 2026}

\begin{document}

\maketitle

\section{Introduction}

\texttt{wmuxd} is a small background daemon that watches your AI coding agent sessions (Claude Code, Codex) for notification events -- ``I need your input,'' ``permission needed,'' etc -- and gives you one place to see and react to them, instead of alt-tabbing between a dozen terminal windows to find out which one is actually waiting on you.

\texttt{wmux} is the CLI you use to talk to it: spawn sessions, open terminal panes, wire up hooks, and watch for notifications live.

For architecture and setup rationale, see \texttt{README.md}. For development history and known gotchas discovered during testing, see \texttt{NOTES.md}.

\section{Before You Start: Where Does Your Agent Actually Run?}

This is the single most important thing to get right, and it's not always obvious. Check:

\begin{lstlisting}
where claude
where codex
\end{lstlisting}

If either prints a \texttt{.exe} path (e.g.\ \texttt{C:\textbackslash Users\textbackslash you\textbackslash .local\textbackslash bin\textbackslash claude.exe}), that agent is a \textbf{native Windows install} -- even if you also happen to have WSL distros on the machine for other things. Only treat an agent as \textbf{WSL-based} if it's actually installed inside a distro:

\begin{lstlisting}
wsl -d <distro> -- which claude
\end{lstlisting}

Everything below is split into ``native Windows'' and ``WSL'' examples -- use the one that matches what you just found. Don't assume; check.

\section{Installation}

\begin{enumerate}
  \item Build or grab the binaries -- \texttt{wmuxd.exe}/\texttt{wmux.exe} for native Windows use, \texttt{bin/linux-amd64/wmuxd}/\texttt{wmux} if you also need a WSL-resident daemon.
  \item Put \texttt{wmux.exe}/\texttt{wmuxd.exe} somewhere permanent (e.g.\ \texttt{C:\textbackslash wmux\textbackslash}) and add that folder to your PATH so you can just type \texttt{wmux} from any terminal:
\begin{lstlisting}
$current = [Environment]::GetEnvironmentVariable("Path", "User")
[Environment]::SetEnvironmentVariable("Path", "$current;C:\wmux", "User")
\end{lstlisting}
  (Open a \textbf{new} terminal afterwards -- existing ones won't see the updated PATH.)
  \item If you also need the WSL-resident daemon, copy the Linux binaries somewhere on the distro's PATH:
\begin{lstlisting}
sudo cp bin/linux-amd64/wmux bin/linux-amd64/wmuxd /usr/local/bin/
\end{lstlisting}
\end{enumerate}

\section{Starting the Daemon}

Native Windows:
\begin{lstlisting}
wmuxd.exe
\end{lstlisting}

WSL-resident:
\begin{lstlisting}
nohup wmuxd > /tmp/wmuxd.log 2>&1 &
\end{lstlisting}

Either way, confirm it's up:
\begin{lstlisting}
curl.exe -s http://127.0.0.1:47823/healthz
# -> ok
\end{lstlisting}

You need \texttt{wmuxd} running before anything else in this manual will work.

\textbf{Sessions survive a daemon restart.} \texttt{wmuxd} snapshots session state to \texttt{\textasciitilde/.wmux/state.json} (override with \texttt{wmuxd --state <path>}, or \texttt{--state ""} to disable) after every change, and restores it on startup -- each restored session's process is re-checked so a session that actually died while the daemon was down comes back correctly marked exited, not stuck showing running.

\section{Wiring the Notification Hook}

This is what makes Claude Code push a message into wmux whenever it's waiting on you, instead of you having to notice a quiet terminal.

\subsection{Native Windows Agent}

Edit \texttt{C:\textbackslash Users\textbackslash\textless you\textgreater\textbackslash .claude\textbackslash settings.json}:

\begin{lstlisting}
{
  "hooks": {
    "Notification": [
      {
        "matcher": "",
        "hooks": [{ "type": "command", "command": "C:\\wmux\\wmux.exe hook-claude" }]
      }
    ]
  }
}
\end{lstlisting}

(Use the full path to \texttt{wmux.exe} here rather than relying on PATH -- the process that invokes hooks doesn't always inherit a freshly updated one.)

\subsection{WSL-Based Agent}

Edit the \emph{WSL-side} \texttt{\textasciitilde/.claude/settings.json} (inside the distro):

\begin{lstlisting}
{
  "hooks": {
    "Notification": [
      { "matcher": "", "hooks": [{ "type": "command", "command": "wmux hook-claude" }] }
    ]
  }
}
\end{lstlisting}

\subsection{Codex (Either Side)}

Edit \texttt{\textasciitilde/.codex/config.toml} (root keys must come before any \texttt{[tables]}):

\begin{lstlisting}
notify = ["wmux", "hook-codex", "--session", "my-project"]
\end{lstlisting}

\section{Running an Agent Session}

\subsection{Example: Native Windows Claude Code, in Your Current Terminal}

\begin{lstlisting}
wmux attach --id my-project --cwd D:\path\to\project -- "C:\Users\you\.local\bin\claude.exe"
\end{lstlisting}

Runs \texttt{claude.exe} right there with a real TTY (colors, readline, prompts all work normally) while registering the session with the daemon under the id \texttt{my-project}. When you exit Claude, it deregisters automatically.

\subsection{Example: Native Windows Claude Code, in a New Windows Terminal Pane}

\begin{lstlisting}
wmux pane --native --id my-project --cwd D:\path\to\project ^
  --cmd "C:\Users\you\.local\bin\claude.exe" --split right
\end{lstlisting}

Opens a new side-by-side split running the same thing. \texttt{--split} also takes \texttt{down} (stacked split) or \texttt{tab} (new tab, the default if you omit \texttt{--split}).

\texttt{wt.exe}'s own \texttt{-V}/\texttt{-H} flags name a split after the orientation of the \emph{dividing line}, which is backwards from what most people mean by ``vertical''/``horizontal'' -- verified by screenshot that \texttt{-V} actually produces left/right, not top/bottom. \texttt{wmux pane} uses \texttt{right}/\texttt{down} instead to sidestep that confusion entirely.

\paragraph{Gotcha.} If the exe path has a space in it (a common case -- usernames like \texttt{Peter Kure} do this), don't wrap it in embedded double quotes inside \texttt{--cmd}. PowerShell 5.1 has a known bug marshalling arguments with literal embedded \texttt{"} to native programs, and it can silently swallow trailing flags like \texttt{--split}. Instead, use the path's 8.3 short form, which has no spaces and needs no quoting:

\begin{lstlisting}
(New-Object -ComObject Scripting.FileSystemObject).GetFile(
  "C:\Users\you\.local\bin\claude.exe").ShortPath
# -> C:\Users\PETERK~1\LOCAL~1\bin\claude.exe
\end{lstlisting}

\begin{lstlisting}
wmux pane --native --id my-project --cwd D:\path\to\project ^
  --cmd "C:\Users\PETERK~1\LOCAL~1\bin\claude.exe" --split right
\end{lstlisting}

\subsection{Example: WSL-Based Codex, Headless (Batch Run, No TTY)}

For scripted/background runs where you don't need to type anything:

\begin{lstlisting}
wmux new --id nightly-refactor --cwd /home/you/my-project \
  --cmd "codex exec 'run the migration'"
\end{lstlisting}

\texttt{--distro} is optional here -- if omitted, \texttt{wsl.exe} uses your system's actual default distro. Only pass it if you want a non-default one:

\begin{lstlisting}
wmux new --id nightly-refactor --cwd /home/you/my-project \
  --cmd "codex exec '...'" --distro Ubuntu-22.04
\end{lstlisting}

\subsection{Example: WSL-Based Claude Code, Interactive}

\begin{lstlisting}
wmux attach --id my-project --cwd /home/you/my-project -- claude
\end{lstlisting}

\subsection{Example: WSL-Based Claude Code, New Windows Terminal Pane}

\begin{lstlisting}
wmux pane --id my-project --cwd /home/you/my-project --cmd claude --split right
\end{lstlisting}

Note: no \texttt{--native} here -- that flag is specifically for the Windows path. Plain \texttt{wmux pane} always launches inside WSL via \texttt{wsl.exe}.

\texttt{--cmd} can be a compound command (semicolons, pipes) and it'll still work correctly -- the underlying quoting through \texttt{wt.exe} is handled for you:

\begin{lstlisting}
wmux pane --id build-and-run --cwd /home/you/my-project \
  --cmd "npm install; npm run dev" --split down
\end{lstlisting}

\section{Watching for Notifications}

Two ways to check in on things.

\subsection{Snapshot}

Current state of every known session (git branch, listening ports, last notification, running/exited):

\begin{lstlisting}
wmux list
\end{lstlisting}

Example output:

\begin{lstlisting}
my-project             running    /home/you/my-project   branch=main   ports=[3000] note="Claude is waiting for your input"
nightly-refactor       exited     /home/you/my-project   branch=main   ports=[]     note=""
\end{lstlisting}

\texttt{ports} is scoped to the session's own process tree, not every listening port on the machine -- with one exception: a WSL-targeted \texttt{wmux new}/plain \texttt{wmux pane} session on a \textbf{Windows-native} daemon falls back to listing every port inside the distro, since that session's tracked PID (the Windows-side \texttt{wsl.exe} frontend) has no correlation to PIDs inside WSL's own namespace.

\subsection{Live Feed}

Streams notifications as they happen, useful to leave running in a spare terminal:

\begin{lstlisting}
wmux watch
\end{lstlisting}

Example output as it happens:

\begin{lstlisting}
watching for notifications... (Ctrl+C to stop)
[14:32:07] my-project: Claude is waiting for your input
[14:35:51] nightly-refactor: Codex finished a turn
\end{lstlisting}

\section{Manual/Test Notifications}

Useful for testing your setup without waiting for a real agent event:

\begin{lstlisting}
wmux notify "test message" --session my-project
\end{lstlisting}

\textbf{Common mistake:} forgetting \texttt{--session}. Without it, the notification still gets pushed (you'll see it in \texttt{wmux watch}), but with an empty session ID -- it won't attach to any session in \texttt{wmux list}, so you won't see it reflected there. Always pass \texttt{--session <id>} matching an id you're tracking with \texttt{wmux new}/\texttt{wmux attach}.

\section{Ending a Session}

Just exit the agent (Ctrl+D, \texttt{/exit}, closing the process) -- \texttt{wmux attach} notices and deregisters automatically, including on a non-zero exit code.

To end it remotely instead of from inside the session:

\begin{lstlisting}
wmux close --id my-project
\end{lstlisting}

Kills the session's tracked process -- the daemon-owned process for \texttt{wmux new}, or the registered PID for \texttt{wmux attach}/\texttt{wmux pane}. Verified: this kills the real OS process (confirmed via process list, not just daemon state) and deregisters the session immediately.

\textbf{If you opened it via \texttt{wmux pane}:} ending the process (whether by exiting the agent yourself or via \texttt{wmux close}) also removes the pane itself from Windows Terminal's layout -- \texttt{wmux pane} opens every pane on an auto-installed \texttt{wmux} WT profile with \texttt{closeOnExit: "always"}, whose fixed commandline (\texttt{wmux pane-exec}) owns the pane's whole process chain. (Panes opened by pre-profile wmux versions passed the commandline straight through \texttt{wt.exe}, which never honors \texttt{closeOnExit} -- those still linger as inert panes and can only be closed by hand.)

\section{Switching Focus}

\begin{lstlisting}
wmux focus --id my-project     # focus that session's pane/tab
wmux focus --dir right         # move focus one pane right
\end{lstlisting}

Run from the Windows side, like \texttt{wmux pane}. \texttt{--id} finds the pane by its title -- every \texttt{wmux pane} keeps the session ID as its fixed pane title -- via UI Automation: it foregrounds the right Windows Terminal window, selects the tab, and puts keyboard focus on the exact pane, including one half of a split. \texttt{--dir} is relative movement (\texttt{left}/\texttt{right}/\texttt{up}/\texttt{down}) within the most recently used WT window -- handy for an agent that just opened a pane next to itself.

\section{Troubleshooting}

\begin{tabular}{|p{4cm}|p{4.5cm}|p{6cm}|}
\hline
\textbf{Symptom} & \textbf{Likely cause} & \textbf{Fix} \\
\hline
Session exits instantly, no output (WSL commands) &
Bad or missing \texttt{--distro} &
\texttt{wsl -l -v} to see real distro names; omit \texttt{--distro} to use your default, or pass the correct name \\
\hline
\texttt{wmux hook-claude}/\texttt{hook-codex} says ``could not reach wmuxd'' &
Hook wired to the wrong side &
A native Windows agent's hook needs a native \texttt{wmuxd}; a WSL agent's hook needs a WSL-resident \texttt{wmuxd} -- they're on separate network namespaces unless WSL2 mirrored networking is on \\
\hline
\texttt{wmux pane} opens a window but the session never shows in \texttt{wmux list} &
\texttt{wmux} not on PATH inside the target WSL distro, or \texttt{wmuxd} isn't running there &
\texttt{wsl -d <distro> -- which wmux}; start \texttt{wmuxd} inside the distro \\
\hline
\texttt{wmux pane --native --cmd "..."} seems to drop a trailing flag like \texttt{--split} &
PowerShell 5.1 embedded-quote bug (see the gotcha above) &
Use the 8.3 short path instead of a quoted long path \\
\hline
\texttt{wmux notify ...} doesn't show up against the right session in \texttt{wmux list} &
Forgot \texttt{--session} &
Add \texttt{--session <id>} \\
\hline
\texttt{wmux close --id X} succeeds but the pane/tab is still visible &
Pane opened by a pre-profile wmux version (commandline passed through \texttt{wt.exe} never honors \texttt{closeOnExit}) &
Close that pane by hand; panes opened by current \texttt{wmux pane} close themselves \\
\hline
\texttt{wmux pane} opens a plain shell instead of the agent &
The \texttt{wmux} WT profile fragment wasn't imported yet (first-ever run) &
Retry \texttt{wmux pane}; if it persists, restart Windows Terminal once \\
\hline
\texttt{wmux focus --id X} says not found &
Pane title no longer matches the session ID (pane opened by an old wmux version, or session not opened via \texttt{wmux pane}) &
Use \texttt{wmux focus --dir} instead, or reopen the pane with current \texttt{wmux pane} \\
\hline
\end{tabular}

\end{document}
